\chapter{\MakeUppercase{Conclusion and Further enhancements}}

\section{Conclusion}
The Badal project successfully demonstrates the feasibility and benefits of an automated, AI-driven cloud security auditing platform. By integrating resource discovery, dependency mapping, and vulnerability correlation into a single cohesive system, Badal provides cloud administrators with a much-needed "security-first" view of their infrastructure.

The project highlights three key findings:
\begin{enumerate}
    \item \textbf{Visualization is Critical:} Representing cloud infrastructure as a dependency graph makes complex security risks immediately apparent.
    \item \textbf{AI Bridging the Gap:} Large Language Models are highly effective at transforming raw security data into actionable remediation steps.
    \item \textbf{Integrated Risk Scoring:} A multi-factor risk score is more effective for prioritization than isolated vulnerability scores.
\end{enumerate}

Badal offers a significant improvement over manual auditing processes, reducing the time required to identify and solve security issues while increasing the overall security posture of the cloud environment.

\section{Further Enhancements}
While the current version of Badal provides a strong foundation, several areas for future improvement have been identified:

\subsection{Multi-Cloud Support}
The current architecture is modular, but scanner implementations are primarily for AWS. Future versions could include scanners for Google Cloud Platform (GCP) and Microsoft Azure to provide a unified multi-cloud security view.

\subsection{Real-time Monitoring}
Implementing a real-time event-driven scanner using AWS CloudTrail and EventBridge would allow Badal to respond to security events as they happen, rather than relying on manual on-demand scans.

\subsection{Automatic Remediation}
For a subset of well-defined, low-risk security issues (e.g., enabling encryption on an S3 bucket), Badal could offer a "One-Click Fix" feature that automatically applies the suggested remediation.

\subsection{Predictive Analysis}
By analyzing historical scan data, the system could potentially predict which resources are most likely to develop vulnerabilities or configuration drifts, allowing for proactive security hardening.

\subsection{Enhanced Policy Analysis}
Integrating a dedicated IAM policy simulator would allow for deeper analysis of permissions and help in achieving a true "Least Privilege" architecture.

% add contibutions after conclusion and further enhancements
\section*{Contributions}
The development of Badal involved contributions across several domains:
\begin{itemize}
    \item \textbf{Infrastructure Research:} Identifying critical cloud resource patterns and security pitfalls.
    \item \textbf{Backend Development:} Implementing the FastAPI server and Boto3 scanning logic.
    \item \textbf{Algorithm Design:} Developing the dependency mapping logic and the multi-factor risk scoring system.
    \item \textbf{AI Integration:} Crafting and testing LLM prompts for accurate remediation generation.
    \item \textbf{System Validation:} Creating test environments and verifying security findings.
\end{itemize}
